\documentstyle[%


righttag%


,11pt%


,amssymb%


,russian%               remove in Eng.


,izvru%                 Set izven% in Eng.


]{amsart}





%





\newcommand{\const}{\operatorname{const}}


\newcommand{\tts}[1]{}





\theoremstyle{plain}


\theorembodyfont{\it}


\newtheorem{thmnonum}{\hskip\parindent Теорема}


\newtheorem{lemnonum}{\hskip\parindent Лемма}


\renewcommand{\thethmnonum}{}


\renewcommand{\thelemnonum}{}


\renewcommand{\baselinestretch}{0.98}








%\hoffset=-15mm%                  For Eng.


\setcounter{page}{41}


\god{2004}


\nomer{3~(502)} %                        Remove in Eng.


%\nomer{Vol.\,48, No.\,3}%               Set in Eng.


%\str{\pageref{begin}--\pageref{end}}%  Set in Eng.


\udk{517.43}





\author{\it В.Б.\,Левенштам}


% NB:   ^^ remove in Eng.


\title{Построение старших приближений метода усреднения для


параболических начально-краевых задач методом пограничного слоя}


\thanks{Работа выполнена при финансовой поддержке


Российского фонда фундаментальных исследований (грант


\No\,01-01-00678) и научной программы ``Университеты России'' (грант


\No\,УР.04.01.029).}





\date{}


%\pagestyle{myheadings}%         Set in Eng.


\pagestyle{plain} %              Remove in Eng.


\begin{document}


%\thispagestyle{plain}%          Set in Eng.


\maketitle


\label{begin}











Первые работы, посвященные построению с математическим


обоснованием старших приближений метода усреднения для


полулинейных параболических начально-краевых задач произвольного


порядка, принадлежат И.Б.\,Симоненко [1], [2]. В них автор,


отправляясь от решения соответствующей усредненной задачи и


используя метод Ньютона--Канторовича и некоторые идеи метода


усреднения [3], строит рекуррентную последовательность линейных


параболических задач, решения которых являются старшими


приближениями решения исходной (возмущенной) задачи. Правые же


части этих линейных задач зависят от большого


параметра --- частоты осцилляций коэффициентов возмущенной задачи,


так что требуется еще асимптотическое интегрирование этих линейных


задач, которое, как и полагал И.Б.\,Симоненко, можно осуществить


методом пограничного слоя [4] с соответствующим обоснованием. В


данной статье построение старших приближений для первой


параболической начально-краевой задачи произвольного порядка


осуществляется с обоснованием, по-видимому, более естественно ---


путем применения метода пограничного слоя непосредственно к


исходной полулинейной задаче. Старшие приближения аппроксимируют


решение возмущенной задачи по сколь угодно сильной г\"ельдеровой


норме. При этом, однако, на осциллирующие с нулевым средним


слагаемые исходного уравнения пришлось наложить существенно более


жесткие ограничения (см. условие (\ref{lev:4}) ниже), нежели те,


которые необходимы для принадлежности


этого решения соответствующему г\"ельдерову пространству. Для первой


параболической начально-краевой задачи второго порядка


используемый здесь алгоритм построения старших приближений


обоснован в [5].


\medskip





$1^\circ$. Пусть $m$, $k$ и $s$ --- натуральные числа, $T>0$,


$\Omega$ --- ограниченная область евклидова пространства $\Bbb R^m$


с $C^\infty$-гладкой границей $\partial \Omega$. В цилиндре


$Q=\overline{\Omega}\times[0,T]$ с боковой поверхностью


$\Gamma=\partial\Omega\times(0,T]$ рассмотрим в обычной


классической постановке задачу


\begin{gather}\label{lev:1}


\frac{\partial u}{\partial t}-\sum_{|\alpha|=2k}


a_\alpha(x)D^\alpha u=\sum_{|\ell|\leqslant s} f_\ell(x,t,\delta^{2k-1}u)


e^{i\ell\omega t},\\


%


\label{lev:2}


u\big|_\Gamma=\frac{\partial u}{\partial


n}\Big|_\Gamma=\dotsb=\frac{\partial^{k-1} u}{\partial


n^{k-1}}\Big|_\Gamma=0,\quad


u(x,0)=\varphi(x).


\end{gather}


Здесь $\alpha=(\alpha_1,\dotsc,\alpha_m)$, $D^\alpha


u\equiv\frac{\partial^{|\alpha|}u}


{\partial x_1^{\alpha_1}\partial x_2^{\alpha_2}\dotsc


\partial x_m^{\alpha_m}}$,


$x=(x_1,\dotsc,x_m)$,


$\delta^{2k-1}u$ --- вектор--функция, составленная из функции $u$


и ее всевозможных производных по $x$ до порядка $2k-1$


включительно, $n$ --- внутренняя нормаль к $\Gamma$, $\omega$ ---


большой параметр. Относительно функций $a_\alpha$, $f_\ell$ и


$\varphi$ предположим следующее.





1. Вещественные функции $a_\alpha(x)$ бесконечно дифференцируемы


при $x\in \overline{\Omega}$, и для задачи\linebreak


(\ref{lev:1})--(\ref{lev:2}) выполнено условие равномерной


параболичности, т.\,е. при всех $x\in \overline{\Omega}$ и


ненулевых векторах $\sigma\in\Bbb R^m$


$$


(-1)^{k+1} \sum_{|\alpha|=2k} a_\alpha(x)\sigma^\alpha


\geqslant c|\sigma|^{2k},\quad c=\const>0.


$$





2. Обозначим через $p$ размерность вектора $(\delta^{2k-1}u)(x)$.


Пусть $D$ --- ограниченная область комплексного пространства $\Bbb


C^p$, а $N_0$ --- некоторое целое неотрицательное число,


конкретнее о котором будет сказано в теореме. Предположим, что


$f_\ell(x,t,e)$, где $e=(e_1,\dotsc,e_p)\in D$,


$|\ell|\leqslant s$ --- комплекснозначные функции, которые


определены на множестве $Q\times D$ и бесконечно дифференцируемы,


причем при $0<|\ell|\leqslant s$ справедливы равенства


\begin{equation}\label{lev:4}


\frac{\partial^j}{\partial t^j} D^\alpha f_\ell(x,t,e)


\Big|_{(x,t)\in \partial\Omega\times\{0\}}=0,\quad


j+|\alpha|\leqslant N_0.


\end{equation}





Наряду с задачей (\ref{lev:1})--(\ref{lev:2}) рассмотрим в $Q$


усредненную задачу


\begin{gather}\label{lev:5}


\frac{\partial v}{\partial t}-\sum_{|\alpha|=2k}


a_\alpha(x)D^\alpha v =f_0(x,t,\delta^{2k-1}v),\\


%


\label{lev:6}


v\big|_\Gamma=\frac{\partial v}{\partial


n}\Big|_\Gamma=\dotsb=\frac{\partial^{k-1} v}{\partial


n^{k-1}}\Big|_\Gamma=0,\quad v(x,0)=\varphi(x).


\end{gather}





3.  Предположим, что задача (\ref{lev:5})--(\ref{lev:6}) имеет


бесконечно дифференцируемое по переменным\linebreak $x$, $t$ решение


$v(x,t)=u_0(x,t)$ такое, что при $(x,t)\in Q$ имеет место


$(\delta^{2k-1}u_0)(x,t)\in D$.





Согласно [6] при больших $\omega$ существует решение $u_\omega$


задачи (\ref{lev:1})--(\ref{lev:2}) и для некоторого $\gamma_0\in


(0,1)$ справедливо соотношение


$$


%%\begin{equation}\label{lev:8}


\lim_{\omega\to\infty} \sum_{|\alpha|\leqslant 2k-1}


\|D^\alpha[u_\omega-u_0]\|_{C^{\gamma_0,\gamma_0/2k}}=0.


%%\end{equation}


$$


Здесь через $C^{h,h/2k}\equiv C_{x,t}^{h,h/2k}(Q)$, $h\geqslant 0$,


обозначено г\"ельдерово пространство функций $u(x,t)$, заданных в


цилиндре $Q$, которое для нецелых $h$ определено в [7], а для


целых $h$ доопределяется естественным образом. Функция $u_0$


называется первым приближением решения $u_\omega$ по указанной


норме.


\medskip





$2^\circ$. Перейдем к построению старших приближений метода


усреднения в сильных г\"ельдеровых нормах. Для этого предварительно


введем в замыкании $\overline{\Omega}_\eta$ пограничной полоски


$\Omega_\eta$ толщины $\eta>0$ криволинейную систему координат


$(\varphi,r)=(\varphi_1,\dotsc,\varphi_{m-1},r)$ следующим образом.


Определим отображение $\partial\Omega\times[0,\eta]\to


\overline{\Omega}_\eta$ по закону $(\varphi,r)\to \varphi+n_\varphi


r$, где $\varphi$ --- точка на $\partial\Omega$, имеющая местные


ортогональные координаты $\varphi$, а $n_\varphi$ --- вектор


внутренней нормали к $\partial\Omega$ в точке $\varphi$. Число


$\eta$ берется столь малым, что указанное отображение обратимо.





Старшие приближения решения $u_\omega$ задачи


(\ref{lev:1})--(\ref{lev:2}) ищем в виде


\begin{equation}


\overset{n}{u}{}(x,t,\omega)=\sum_{j=0}^n \omega^{-j/2k}


u_j(x,t)+ \omega^{-1}\sum_{j=0}^n \omega^{-j/2k}


[v_j(x,t,\omega t)+


w_j(\varphi,\rho,t)+z_j(\varphi,\rho,t,\omega t)],


\label{lev:9}


\end{equation}


где $\rho=r\omega^{1/2k}$, $u_j$ и $v_j$ --- регулярные [4], а


$w_j$ и $z_j$ --- погранслойные [4] функции, причем функции


$v_j(x,t,\tau)$ и $z_j(\varphi,\rho,t,\tau)$ ---


$2\pi$-периодические по $\tau$ с нулевым средним


$$


\langle v_j \rangle(x,t)=\frac{1}{2\pi}


\int_0^{2\pi} v_j(x,t,\tau)\,d\tau=0,\quad


\langle z_j\rangle=0,


$$


а функции $w_j$ и $z_j$  в области $\Omega\setminus


\Omega_\eta$ равны нулю.





В пункте $3^\circ$ будет описан алгоритм построения коэффициентов


приближений $\overset{n}{u}{}$, при котором справедлива





\begin{thmnonum}


Существует такое число $\omega_0>0$, что при $\omega>\omega_0$ 


для любого $h\geqslant 0$ и любого натурального $n$ найдется


такое число $N_0$, что при выполнении условий п.\,$1^\circ$ с этим


значением $N_0$ в $(\ref{lev:4})$ имеют место оценки


$$


\|u_\omega-\overset{n}{u}{}\|_{C^{h,h/2k}(Q)}\leqslant


C_{n,h} \omega^{-\frac{n+1-\max(h-2k,0)}{2k}},\quad


C_{n,h}=\const.


$$


Построение приближения $\overset{n}{u}{}$ при известном решении


$u_0$ усредненной задачи сводится к решению~$n$


однозначно разрешимых линейных первых


начально-краевых задач для параболических уравнений 


$Pu=\varphi(x,t)$ в $Q$ порядка $2k$ с единым дифференциальным выражением $P$.


\end{thmnonum}





\begin{note}


Предположения п.\,$1^\circ$


о бесконечной дифференцируемости границы $\partial\Omega$ и


функций $a_\alpha$, $f_\ell$ и $u_0$, естественно, могут быть


заменены на требования их непрерывной дифференцируемости вплоть до


определенных конечных порядков (зависящих от $h$ и $n$) так, что


при этой замене утверждения теоремы сохраняются.


\end{note}





\begin{note}


В условии (\ref{lev:4}) функцию


$f_\ell(x,t,e)$ можно заменить на функцию


$\widetilde{f}_\ell(x,t)=f_\ell[x,t,\delta^{2k-1}u_0(x,t)]$.


\end{note}





$3^\circ$. Опишем схему нахождения неизвестных функций $u_j$,


$v_j$, $w_j$ и $z_j$, входящих в представление (\ref{lev:9})


старших приближений. При этом будем пользоваться, как принято в


методе пограничного слоя [4], представлением эллиптической части


уравнения (\ref{lev:1}) в переменных $(\varphi,\rho)$:


\begin{equation}


\sum_{|\alpha|=2k} a_\alpha(x) D^\alpha u=


\omega\bigg[ b_{2k}(\varphi) \frac{\partial^{2k}u}{\partial \rho^{2k}}


+\sum_{j=1}^N \omega^{-j/2k} M_j(\varphi,\rho)u+


\omega^{-\frac{N+1}{2k}}M_{N+1}(\varphi,\rho;r)u \bigg],


\label{lev:10}


\end{equation}


где $b_{2k}\in C^\infty(\partial\Omega)$, $(-1)^{k+1}b_{2k}>0$; $N$ ---


натуральное число; $M_j$, $1\leqslant j\leqslant N$, и


$M_{N+1}$ --- дифференциальные выражения относительно $\varphi$,


$\rho$ с бесконечно дифференцируемыми по $(\varphi,\rho)$ и


соответственно по $(\varphi,r)$ коэффициентами.





Подставим выражение (\ref{lev:9}) приближения $\overset{n}{u}{}$


вместо $u$ в равенства (\ref{lev:1})--(\ref{lev:2}), разложим


функции $f_\ell(x,t,\delta^{2k-1}\overset{n}{u}){}$ в ряды Тейлора


по последним $p$ аргументам с центром $\delta^{2k-1}u_0$ и,


учитывая (\ref{lev:10}), приравняем коэффициенты при одинаковыx


степенях $\omega$, стоящих в разных частях каждого из полученных


равенств, отдельно для регулярных и погранслойных функций. Каждое


из составленных таким образом уравнений разобьем на уравнения для


плавных, т.\,е. не зависящих от $\tau=\omega t$ и осциллирующих


по $\tau$ с нулевым средним коэффициентов. В результате для $u_0$


получаем полулинейную усредненную задачу


(\ref{lev:5})--(\ref{lev:6}), которая по условию разрешима. Для


функции $v_0(x,t,\tau)$ получаем задачу


$$


\frac{\partial v_0}{\partial\tau}=\sum_{0<|\ell|\leqslant s}


f_\ell(x,t,\delta^{2k-1}u_0)e^{i\ell\tau},\quad


\langle v_0\rangle=0,


$$


где $x$, $t$ играют роль параметров, и которая, очевидно,


однозначно разрешима. При этом в силу условия (\ref{lev:4})


справедливы соотношения


\begin{equation}\label{lev:11}


\frac{\partial^j}{\partial t^j}D^\alpha v_0(x,t,\tau)


\Big|_{(x,t)\in \partial\Omega\times\{0\}}=0,\quad


j+|\alpha|\leqslant N_0.


\end{equation}





Функции $u_i$, $v_i$, $i\geqslant 1$, являются решениями линейных


задач вида


\begin{gather}\label{lev:12}


\displaybreak[3]


\frac{\partial u_i}{\partial t}-\sum_{|\alpha|=2k}


a_\alpha(x)D^\alpha u_i-


[(D_e f_0)(x,t,\delta^{2k-1}u_0)](\delta^{2k-1}u_i)=\psi_i(x,t),\\


%


u_i\big|_\Gamma=-w_{i-2k}\Big|_{\rho=0},\ \


\frac{\partial u_i}{\partial n}\Big|_\Gamma=


-\frac{\partial w_{i-2k+1}}{\partial \rho}\Big|_{\rho=0}, \dotsc,


\frac{\partial^{k-1} u_i}{\partial n^{k-1}}\Big|_\Gamma=\hspace*{50mm}\notag\\


%


\label{lev:13}


\hspace*{50mm}=-\frac{\partial^{k-1} w_{i-k-1}}{\partial


\rho^{k-1}}\Big|_{\rho=0},\ \ u_i(x,0)=-v_{i-2k}(x,0,0),\\


%%


\label{lev:15}


\frac{\partial v_i(x,t,\tau)}{\partial \tau}=\chi_i(x,t,\tau),\quad


\langle v_i\rangle=0.


\end{gather}


Здесь $[(D_e f_0)(x,t,\overset{\circ}{e}{})](e)=\sum\limits_{i=1}^ p


\frac{\partial f_0(x,t,\overset{\circ}{e}{})}{\partial e_i}e_i$;


$\psi_i$ и $\chi_i$ --- бесконечно дифференцируемые функции,


которые выражаются через $u_j$, $v_j$, $j\leqslant i-1$, а $v_i$,


$w_i$ при $i<0$ полагаются нулями.





Задачи для функций $w_i$, $z_i$, $i\geqslant 0$, в погранполоске


имеют вид


\begin{gather}\label{lev:16}


b_{2k}(\varphi)\frac{\partial^{2k}w_i}{\partial \rho^{2k}}=


\lambda_i(\varphi,\rho,t),\quad


w_i\big|_{\rho=\infty}=0,\\


%


\label{lev:17}


\frac{\partial z_i(\varphi,p,t,\tau)}{\partial \tau}-b_{2k}(\varphi)


\frac{\partial^{2k} z_i(\varphi,p,t,\tau)}{\partial


\rho^{2k}}=\mu_i(\varphi,\rho,t,\tau),\quad


z_i\big|_{\rho=\infty}=0,\\


%


\label{lev:18}


z_i\big|_{\rho=0}=-v_i\big|_\Gamma,\ \


\frac{\partial z_i}{\partial \rho}\Big|_{\rho=0}=


-\frac{\partial v_{i-1}}{\partial


n}\Big|_\Gamma, \dotsc ,


\frac{\partial^{k-1} z_i}{\partial \rho^{k-1}}\Big|_{\rho=0}=


-\frac{\partial v_{i-k+1}}{\partial


n^{k-1}}\Big|_\Gamma,


\end{gather}


где $\lambda_i$, $\mu_i$ --- бесконечно дифференцируемые функции,


которые выражаются через $u_j$, $v_j$, $w_j$, $z_j$, $j\leqslant


i-1$.





Докажем теперь, что задачи (\ref{lev:12})--(\ref{lev:13}),


(\ref{lev:15}), (\ref{lev:16}) и (\ref{lev:17})--(\ref{lev:18})


однозначно разрешимы. Пусть для простоты в условии


(\ref{lev:4}) \ $N_0=\infty$. Итак, определив функции $u_0$, $v_0$, 


переходим к задаче (\ref{lev:16}) при $i=0$


$$


b_{2k}(\varphi)\frac{\partial^{2k}w_0}{\partial


\rho^{2k}}=0,\quad   w_0\big|_{\rho=\infty}=0


$$


с единственным решением $w_0=0$ и к задаче


(\ref{lev:17})--(\ref{lev:18}) при $i=0$


\begin{gather*}


%% \label{lev:19}


\frac{\partial z_0}{\partial \tau}-b_{2k}(\varphi)


\frac{\partial^{2k} z_0}{\partial


\rho^{2k}}=0,\quad   z_0\big|_{\rho=\infty}=0,\\


%


%% \label{lev:20}


z_0\big|_{\rho=0}=-v_0\big|_{\partial\Omega},\quad


\frac{\partial z_0}{\partial \rho}\Big|_{\rho=0}=0, \dotsc,


\frac{\partial^{k-1} z_0}{\partial


\rho^{k-1}}\Big|_{\rho=0}=0,


\end{gather*}


которая в силу соотношений (\ref{lev:11}) и сформулированной ниже


леммы также однозначно разрешима и при любых целых неотрицательных


числах $j$ для ее решения $z_0$ справедливы равенства


\begin{equation}\label{lev:21}


\frac{\partial^j z_0(\varphi,\rho,t,\tau)}{\partial t^j}


\Big|_{t=0}=0.


\end{equation}


Задача (\ref{lev:12})--(\ref{lev:13}) при $i=1$ является линейной


однородной начально-краевой задачей с единственным решением


$u_1=0$. Предположим теперь, что при всех $0\leqslant i<i_0$, где


$i_0$ --- некоторое натуральное число, найдены функции $u_{i+1}$,


$v_i$, $w_i$ и $z_i$, причем для $u_{i+1}$ и $v_i$ выполнены те же


соотношения (\ref{lev:11}), что и для $v_0$, а для $w_i$ и


$z_i$ --- те же соотношения (\ref{lev:21}), что и для $z_0$.


Предположим еще, что функции $\lambda_i$ являются конечными


суммами функций вида $c(\varphi,t)\rho^n e^{-\lambda\rho}$, где


$n\geqslant 0$, $\mathrm{Re\,}\lambda>0$, а функции $c(\varphi,t)$


удовлетворяют соотношениям


\begin{equation}\label{lev:22}


\frac{\partial^j c(\varphi,t)}


{\partial t^j}\Big|_{t=0}=0,\quad j=0,1,\dotsc


\end{equation}


Функции же $\mu_i$ являются тригонометрическими полиномами по


$\tau$ с нулевым средним, коэффициенты которых имеют ту же


структуру, что и $\lambda_i$. Докажем, что функции


$u_{i_0+1}$, $v_{i_0}$, $w_{i_0}$ и $z_{i_0}$ однозначно


определены (как решения рассматриваемых задач) и для них


справедливы высказанные предположения. Действительно, в силу этих


предположений данные линейной начально-краевой задачи


(\ref{lev:12})--(\ref{lev:13}) при  $i=i_0+1$ сколь угодно гладкие


и для нее выполнены условия согласования сколь угодно высокого


порядка, откуда следует ее однозначная разрешимость и


справедливость для решения $u_{i_0+1}$ соотношений


$$


\frac{\partial^j}{\partial t^j} D^\alpha u_{i_0+1}


\Big|_{(x,t)\in \partial\Omega\times\{0\}}=0,\quad


j,|\alpha|=0,1,\dotsc


$$


Легко исследуется и задача (\ref{lev:15}) при $i=i_0$. Здесь


дополнительно учитывается, что $\chi_{i_0}$ --- тригонометрический


многочлен по $\tau$ с нулевым средним. Однозначная разрешимость


задач (\ref{lev:16}) и (\ref{lev:17})--(\ref{lev:18}) и


справедливость для их решений высказанных предположений вытекает


из следующего простого утверждения.





\begin{lemnonum}


Пусть $c_i(\varphi,t)$,


$i=0,1,\dotsc,k$, --- сколь угодно гладкие на множестве


$\partial\Omega\times[0,T]$ функции, удовлетворяющие условиям


$(\ref{lev:22})$, $\lambda(\varphi)$ ---  сколь угодно гладкая на


множестве $\partial\Omega$ функция, $\mathrm{Re\,}\lambda<0$,


$n\geqslant 0$, $\ell\in\Bbb R$, $\ell\ne0$. Тогда задача


\begin{gather*}


i\ell y-b_{2k}(\varphi) \frac{\partial^{2k}y}{\partial


\rho^{2k}}=


c_0(\varphi,t)\rho^n e^{\lambda(\varphi)\rho},\quad


y\big|_{\rho=0}=c_1(\varphi,t),\\


%


\frac{\partial y}{\partial \rho}\Big|_{\rho=0}=c_2(\varphi,t),


\dotsc,


\frac{\partial^{k-1}y}{\partial \rho^{k-1}}\Big|_{\rho=0}=c_k(\varphi,t),\quad


y\big|_{\rho=\infty}=0,


\end{gather*}


однозначно разрешима и ее сколь угодно гладкое решение


$y(\varphi,\rho,t)$ удовлетворяет равенствам


$$


\frac{\partial^j y(\varphi,\rho,t)}


{\partial t^j}\Big|_{t=0}=0,\quad j=0,1,\dotsc


$$


\end{lemnonum}





Итак, показали, что приближения $\overset{n}{u}{}$, заданные


соотношением (\ref{lev:9}), могут быть построены для любого номера


$n$. Доказательство указанной в теореме оценки проводится в два


этапа. На первом этапе методами, которые использовались в работе


[6], эта оценка доказывается для малых $h\geqslant 0$. На втором


этапе она выводится для произвольных $h>0$ из оценки, полученной


на первом этапе, с помощью многократного использования


классической априорной оценки [7] для линейной параболической


задачи.





{\bf Примечание.} Описанный применительно к задаче (\ref{lev:1})--(\ref{lev:2})


алгоритм построения старших приближений сохраняется вместе с обоснованием,


если в уравнении (\ref{lev:1}) сомножители $e^{i\ell\omega t}$


заменить сомножителями $e^{i\lambda_{\ell}\omega t}$,


где $\lambda_\ell$, $|\ell|\leqslant s$, --- произвольный набор


действительных чисел.








\begin{thebibliography}{9}





\bibitem{b1}


Симоненко И.Б. {\em Старшие приближения метода осреднения для


абстрактных параболических уравнений} // Матем. сб. -- 1973. --


Т.\,92. -- \No\,4. -- С.\,541--549.





\bibitem{b2}


Симоненко И.Б. {\em Старшие приближения метода осреднения для


параболических уравнений} // ДАН СССР. -- 1973. -- Т.\,213. --


\No\,6. -- С.\,1255--1257.





\bibitem{b3}


Боголюбов Н.Н., Митропольский Ю.А. {\em Асимптотические методы в


теории нелинейных колебаний}. -- М.: Наука, 1974. -- 503~с.





\bibitem{b4}


Вишик М.И., Люстерник Л.А. {\em Регулярное вырождение и пограничный


слой для нелинейных дифференциальных уравнений с малым параметром}


// УМН. -- 1960. -- Т.\,15. -- Вып.\,4. -- \linebreak С.\,3--95.





\bibitem{b5}


Левенштам В.Б. {\em Старшие приближения метода усреднения для


параболических начально-краевых задач с быстро осциллирующими


коэффициентами} // Дифференц. уравнения. -- 2003. -- Т.\,39. --


\No\,10. -- С.\,1395--1403.





\bibitem{b7}


Симоненко И.Б. {\em Обоснование метода осреднения для абстрактных


параболических уравнений} // Матем. сб. -- 1970. -- Т.\,81.


-- \No\,1. -- С.\,53--61.





\bibitem{b6}


Солонников В.А. {\em О краевых задачах для линейных параболических


систем дифференциальных уравнений общего вида} // Тр. Матем. ин-та АН


СССР. -- 1965. -- Т.\,83. -- С.\,3--162.





\end{thebibliography}





\vskip 2cm





\post{\it Ростовский государственный}{\it Поступила}


\post{\it университет}{11.03.2003}





\label{end}


\end{document}


